\section{Determine the load before synthesizing emitters}
\label{sec:ideal-load}

An emitter-independent test isolates the mechanical question: if the required
mean traction could be supplied exactly, would the proposed pair be an
equilibrium, and what would make it form and remain there? This is an ideal
distributed-load problem, not evidence that an acoustic field realizes it.
The following deductions specialize the variational mechanics of
\cref{sec:mechanics,sec:geometric-inverse}; they introduce no empirical
settling law or additional material data.

\subsection{Two full-domain loads with an explicit sign convention}
Let $D$ be the entire horizontal disk, including both non-optical annuli.
Number the three layers from bottom to top by $0,1,2$, and write their two
interfaces as $z=h_j(x,y)$, $j=0,1$. Use upward normals on \emph{both}
interfaces, not the outward normals of the middle lens. Define
$b_j=(\rho_j-\rho_{j+1})g_{\rm grav}$ and let $f_j$ be upward normal mean
traction, in Pa. Its virtual work is $\int_D f_j\delta h_j\dd x\dd y$:
the graph area factor cancels the normal projection of vertical displacement.
For sealed outer layers, each $\int_D h_j\dd x\dd y$ is fixed; the middle
volume follows from their difference. For compatible full targets $h_j^*$,
\begin{equation}
 f_j^*(x,y)=-\sigma_j\nabla\cdot
       \frac{\nabla h_j^*}{\sqrt{1+|\nabla h_j^*|^2}}
       +b_j h_j^*-\lambda_j,\qquad (x,y)\in D.
 \label{eq:ideal-two-loads}
\end{equation}
The two constants are hydrostatic pressure-difference multipliers, not
independently prescribed absolute pressures. Subtracting each disk average
is one reporting gauge; absolute pressure safety remains a separate issue.
Changing the reference height only changes that gauge. For outward lens
tractions the front sign reverses. Neither load may be silently set to zero
outside the optical pupil. Different admissible annuli give different loads,
even with identical optical conjugates. No unique numerical pressure map
exists until the complete geometry, fluids, fill and support are specified.

\subsection{A conditional three-dimensional stability result for ideal loads}
Hold $f_j^*(x,y)$ fixed as a function of horizontal position while the
interfaces move. Define the loaded potential
\begin{equation}
 \mathcal U[h]=\sum_{j=0}^1\int_D\left[
  \sigma_j\sqrt{1+|\nabla h_j|^2}+\tfrac12 b_jh_j^2-f_j^*h_j
 \right]\dd x\dd y.
 \label{eq:ideal-loaded-potential}
\end{equation}
Suppose $\sigma_j>0$, $b_j\geq0$, the boundary heights are pinned, and
admissible graphs satisfy the fixed volumes and positive layer gaps.
The second variation in \cref{eq:graph-second-variation}, summed over both
interfaces, is strictly positive for any nonzero admissible perturbation.
Indeed, its capillary integrand is bounded below by
$\sigma_j|\nabla\eta_j|^2/(1+|\nabla h_j|^2)^{3/2}$.
The pinned boundary excludes constant perturbations. The admissible graph
set is convex, so $h^*$ is its unique stationary point and strict global
minimizer of $\mathcal U$. This statement allows arbitrary $\eta_j(x,y)$,
including every angular order, not just axisymmetric disturbances. Uniform
coercivity in a neighborhood additionally uses a bound on target slopes.
It is not a result for overturning interfaces or a density-inverted stack.

For smooth incompressible Newtonian mean flow with passive stationary walls,
constant surface tensions, no streaming or other bulk forcing, and exactly
these imposed normal tractions, the mechanical balances give
\begin{equation}
 \frac{\dd}{\dd t}\left[
  \sum_{\ell=0}^2\frac12\int_{\Omega_\ell(t)}\rho_\ell|\bm U_\ell|^2\dd V
  +\mathcal U[h]-\mathcal U[h^*]\right]
 =-\sum_{\ell=0}^2\int_{\Omega_\ell(t)}2\mu_\ell
       \bm D(\bm U_\ell):\bm D(\bm U_\ell)\dd V\leq0.
 \label{eq:ideal-load-energy}
\end{equation}
There is no invented diagonal damping matrix in this identity. Pressure
constraint work vanishes by volume conservation; the two interfaces remain
coupled through the fluid velocity and pressure fields. With sufficient
regularity and coercivity it supplies energy stability about the target.
Asymptotic convergence additionally needs dissipation that excludes persistent
motions, global admissible solutions and compactness or decay estimates.
An energy norm does not alone certify maximum height or optical slope errors.
With zero viscosity, the right side vanishes: oscillation is allowed and
settling cannot be claimed. Thus an inviscid animation cannot answer the
user's maintained-formation question even with perfect loading.

The current hypothetical densities $1200,1100,1000\ \mathrm{kg\,m^{-3}}$
give $b_0=b_1=981\ \mathrm{Pa\,m^{-1}}$ at
$g_{\rm grav}=9.81\ \mathrm{m\,s^{-2}}$, satisfying the sign hypothesis.
These are not characterized laboratory liquids. The result does not transfer
automatically to a lens surrounded by one common liquid: then the two density
jumps have opposite signs unless densities match.

\subsection{Formation requires a load history, not just its endpoint}
Applying $f^*$ from a flat initial state defines a genuine forward initial
value problem. It does not prescribe intermediate shapes. For a time-varying
load, the energy balance includes its time-dependent external work; no
monotonic target-error curve follows from a ramp. A full prescribed trajectory
would instead require solving its compatible velocity and pressure fields
and recovering the interfacial traction from their stress jump, including
inertia and viscosity. It cannot generally be obtained by evaluating only
the static curvature formula along an interpolated shape.

In the specifically declared constant-mass reduction of
\cref{sec:reduced-formation}, the analogous inverse dynamic force is
$F_{\rm req}(t)=M\ddot q_d+\nabla E(q_d)$; a viscous resistance may be added
only after it is derived or independently validated. Endpoint balance removes
the acceleration term but says nothing about transient overshoot. This
identity defines a diagnostic load, not an acoustic realization or a
forward-verified trajectory.

\subsection{Verification gates before returning to source design}
First export capillary, gravitational and gauge-fixed total loads on both
full targets, with geometry, volume and pinning checks. Independently solve
the nonlinear prescribed-load equilibrium from a non-target initial guess;
refine both the load representation and the mechanical grid. Next solve
viscous two-interface formation with independently specified material data
and initial conditions, refine spatial and temporal resolution, and test
three-dimensional disturbances. Report both maximum height errors, slopes,
fixed-detector maximum spot radius and lost rays throughout physical time.
Distinguish a conditional ideal-load pass from any acoustic pass.

Only then ask the acoustic inverse to realize the loads and the required
response to perturbations. At fixed emitter commands the traction is
$f[h,a]$, not the fixed $f^*(x,y)$ above; its shape derivative changes the
restoring operator and invalidates the ideal-load Lyapunov proof unless
controlled. Mean tangential forcing, bulk streaming, wall layers and thermal
feedback also require their own balances. The required acoustic specification
is therefore a full-domain equilibrium load \emph{plus} an admissible dynamic
response, not a single pressure-amplitude image on the pupils.
